%% kapa: exam=1 round=29 year=2018 period=1 form=A subject=civil_law subject_ko=민법 seq=12 number=12
%% source: Q-Net 감정평가사 기출문제 (cst00309) artlSeq=5202808 file=extracted/1008/5202808/2018년 제29회 감정평가사 1차시험 기출문제/2018년 제29회 감정평가사 1차시험 기출문제/1교시 A형.pdf
%% conversion: kordoc 4.24.51 -> tools/gen1.py (auto); review=proofread-sonnet; flags=none
\begin{question}{12}
복대리에 관한 설명으로 옳지 않은 것은? (다툼이 있으면 판례에 따름)
\begin{choices}
  \item 법정대리인은 자신의 책임으로 복대리인을 선임할 수 있다.
  \item 임의대리인은 부득이한 사유가 있는 경우, 복대리인을 선임할 수 있다.
  \item 법정대리인이 부득이한 사유로 복대리인을 선임한 경우, 본인에 대하여 그 선임감독에 관한 책임이 있다.
  \item 임의대리인이 본인의 승낙을 얻어 복대리인을 선임한 경우, 본인에 대하여 그 선임감독에 관한 책임이 없다.
  \item 대리인이 대리권 소멸 후 복대리인을 선임하여 그로 하여금 대리행위를 하도록 한 경우, 복대리인의 대리행위에 대하여 표현대리에 관한 규정이 적용될 수 있다.
\end{choices}
\end{question}
\answer{4}
